\documentclass[12pt]{article}
\usepackage[T1]{fontenc}
\usepackage[utf8]{inputenc}
\usepackage{lmodern}
\usepackage{textcomp}
\usepackage{enumitem}
\usepackage{geometry}
\geometry{margin=1in}
\begin{document}
\begin{center}
Answer Key - Psychology - Undergraduate Level Assessment 1 \
\small (Answers are ordered exactly as in the question paper.)
\end{center}

\section*{Multiple Choice}
\begin{enumerate}[label=\arabic*.]
\item \textbf{Answer:} A
\\ \textit{Options:} A) at the end of the list; B) words that were repeated in the list; C) the longest words in the list; D) words related to a specific topic; E) the last ten words he read
\\ \textit{Note:} The correct answer is at the end of the list because the serial position effect posits that individuals tend to remember items presented at the end of a list better than those in the middle, due to the recency effect associated with short-term memory retention.
\item \textbf{Answer:} A
\\ \textit{Options:} A) basal ganglia; B) hippocampus; C) prefrontal cortex; D) hypothalamus; E) amygdala
\\ \textit{Note:} The basal ganglia are thought to play a crucial role in the Cannon-Bard theory of emotion because they are involved in the integration of emotional responses and the coordination of physical reactions, aligning with the theory's assertion that emotions and physiological responses occur simultaneously.
\item \textbf{Answer:} B
\\ \textit{Options:} A) Homeostatic regulation; B) Hierarchy of needs; C) Self-determination theory; D) Intrinsic motivation; E) Maslow's pyramid of desire
\\ \textit{Note:} The correct answer is the "Hierarchy of needs" because Maslow's theory categorizes human motivations into levels, with basic physiological needs being the most critical for survival, followed by safety, love, esteem, and self-actualization.
\item \textbf{Answer:} E
\\ \textit{Options:} A) symmetry; B) complexity; C) novelty; D) similarity; E) proximity
\\ \textit{Note:} The correct answer is proximity because the letters "c," "a," and "r" are physically closer together in the arrangement, leading our perception to group them as the word "car" rather than the other fragmented parts.
\item \textbf{Answer:} A
\\ \textit{Options:} A) Intermittent reinforcement of the conditioned response; B) Habituation to the conditioned stimulus; C) Transformation of conditioned respondent behavior; D) Emergence of a new unconditioned response; E) No change in conditioned respondent behavior
\\ \textit{Note:} The correct answer is based on the concept of extinction in classical conditioning, where presenting the conditioned stimulus (CS) without the unconditioned stimulus (UCS) leads to a gradual weakening of the conditioned response, but intermittent reinforcement can occur if the CS is occasionally paired with the UCS, thereby maintaining some level of response.
\end{enumerate}

\section*{True / False}
\begin{enumerate}[label=\arabic*.]
\setcounter{enumi}{5}
\item \textbf{Answer:} True
\\ \textit{Options:} A) True; B) False
\\ \textit{Note:} In group counseling, confidentiality is largely upheld through the members' mutual commitment to honor and respect each other's privacy, as formal enforcement mechanisms are limited in such settings.
\item \textbf{Answer:} True
\\ \textit{Options:} A) True; B) False
\\ \textit{Note:} The Draw-a-Person Test evaluates fine motor coordination by analyzing the details and complexity of the drawings produced, which reflect the individual's ability to control their hand movements and spatial awareness.
\item \textbf{Answer:} False
\\ \textit{Options:} A) True; B) False
\\ \textit{Note:} The statement is false because recognition specifically refers to correctly identifying stimuli based on prior experience, while Antonia's mislabeling of the rabbit as a cat indicates a failure in this process, demonstrating a lack of accurate recognition.
\item \textbf{Answer:} True
\\ \textit{Options:} A) True; B) False
\\ \textit{Note:} The statement is true because the presence of other bystanders often leads to the bystander effect, where individuals are less likely to help in an emergency due to diffusion of responsibility and social influence.
\item \textbf{Answer:} True
\\ \textit{Options:} A) True; B) False
\\ \textit{Note:} The statement is true because the phenomenon of afterimages occurs due to the way our visual system processes colors, leading to the perception of complementary colors (in this case, orange) after prolonged exposure to a specific color (yellow).
\end{enumerate}

\section*{Long Form Response}
\begin{enumerate}[label=\arabic*.]
\setcounter{enumi}{10}
\item \textbf{Answer:} The psychological impact of breastfeeding on the parent-child relationship is significant, particularly influenced by parental attitudes during feeding times. When parents approach breastfeeding with positive attitudes-such as warmth, responsiveness, and nurturing-they foster a secure attachment, which is crucial for the child's emotional and social development. Conversely, negative or indifferent attitudes can lead to feelings of insecurity and anxiety in the child, potentially hindering their ability to form healthy relationships later in life. This interaction during feeding not only fulfills the child's nutritional needs but also serves as a critical bonding experience that shapes their sense of safety and trust in their caregivers. Ultimately, the emotional environment created during feeding has lasting effects on the child's development and their perception of relationships.
\\ \textit{Note:} correct because it identifies the crucial role of parental attitudes during breastfeeding in shaping secure attachment and emotional development, emphasizing that positive interactions promote security while negative attitudes can lead to insecurity in the child.
\item \textbf{Answer:} Nonprofit organizations often utilize compliance strategies such as offering free gifts to solicit contributions, leveraging the principle of reciprocity in social psychology. This strategy is effective because it creates a sense of obligation in potential donors; when individuals receive something for free, they may feel compelled to return the favor by making a donation. The psychological underpinning lies in the social norm that encourages giving back when one has benefited, thus enhancing the likelihood of compliance with the organization's request for support. Furthermore, this tactic can increase donor engagement, as the initial gift may serve as a bridge to foster a longer-term relationship between the nonprofit and the donor. Overall, this compliance strategy illustrates how understanding psychological principles can enhance fundraising efforts.
\\ \textit{Note:} correct because it effectively illustrates how the principle of reciprocity in social psychology drives compliance in nonprofit organizations by creating a perceived obligation to reciprocate after receiving a free gift, thereby increasing the likelihood of donations.
\end{enumerate}

\vfill
\noindent\textit{End of Answer Key}
\end{document}